%% F05 --- Functions and graphs
%%
%% Every number in this program that the reader cannot do in their head comes
%% from code/f05_functions.py and is pulled in with \val{}. Arithmetic the
%% program is TEACHING -- x^2 at 0, 1, 2, 3; the shifts; 2x+1 at 0, 1, 2;
%% sigma(0) = 1/2 -- is written inline, because the reader is meant to do it.
%%
%% THE HEADLINE TRAP is section 3's: f(x - 3) moves the graph to the RIGHT.
%% It is elicited at frame 15 under a speed instruction and corrected at 16,
%% and the correction has to explain the reasoning that produced the error --
%% outside the bracket a minus really does move things down, and the reader
%% carried a true rule across the bracket. That is the same reasoning
%% Programs F1, F2, F3 and F4 each caught under a different sign, and the
%% trapbox names those places rather than counting them.
%%
%% THE PAYOFF is section 7's, and it is the reason this program exists in the
%% Foundation part rather than nowhere: a strictly increasing function keeps
%% the order of a list, so it cannot move the argmax. Temperature therefore
%% cannot change which token is most likely -- only how often the others get
%% drawn instead. Both halves are computed in code/f05_functions.py over one
%% fixed logit vector and the invariance is ASSERTED there, so the table on
%% the page cannot quietly stop being true.
%%
%% WHAT THE PAGE MAY NOT SAY, and the script's docstring says why: NOT "the
%% probabilities do not change" -- they change enormously, and the striking
%% figure is the tail rather than the top. Write the invariance of the ARGMAX.
%%
%% Twin: programs/pl/F05-functions-graphs.tex. Frame for frame, macro for
%% macro, number for number; tools/parity.py compares them.

\program{Functions and graphs}
\label{prog:F05}
\index{function!as a machine}
\index{graph!transformations of}

\begin{outcomes}
\outcome{1--7}{say whether a rule is a function of $x$, give the values of $x$
  it accepts, and tell the machine $f$ apart from the number $f(x)$}
\outcome{8--13}{read a value, a crossing point and a flat stretch off a graph,
  and name an interval on which a function increases}
\outcome{14--23}{shift, stretch and reflect a graph, and work out which way
  each move goes from where its number is written}
\outcome{24--30}{say what the weight and the bias each do to a curve, and work
  out where a logistic crosses a half}
\outcome{31--35}{read a composition from the inside out, and say when two
  functions may be wired together at all}
\outcome{36--40}{write down the inverse of a short function, say when one
  exists, and draw its graph from the original}
\outcome{41--47}{say what a strictly increasing function does to a list of
  scores, and use that to decide whether a given change can move the largest
  one}
\end{outcomes}

\begin{quiz}
\item Is $y = x^{2}$ a function of $x$? Is $y^{2} = x$? \teachesat{2--4}
  \answerto{The first is; the second is not. At $x = 4$ the second offers both
  $y = 2$ and $y = -2$, and a function returns one value.}
\item For which values of $x$ does $f(x) = \frac{1}{x - 2}$ give an answer?
  \teachesat{4--5}
  \answerto{Every $x$ except $2$. That is its domain, and the excluded point
  is the one that would divide by zero.}
\item In code, what is the difference between \code{act = relu} and
  \code{act = relu(x)}? \teachesat{6--7}
  \answerto{The first stores the machine, the second stores one number the
  machine has already produced. The parentheses are the difference between
  having a function and having a value.}
\item Sketch $y = 2x + 1$ and say where it crosses the horizontal axis.
  \teachesat{9--10}
  \answerto{At $x = -\num{0.5}$, where $2x + 1 = 0$.}
\item Over which values of $x$ is $y = x^{2}$ increasing? \teachesat{10--11}
  \answerto{For $x \ge 0$. It decreases for $x \le 0$, so \enquote{increasing}
  is a statement about an interval and not about the function as a whole.}
\item The graph of $y = f(x)$ is given. Which way, and how far, does the graph
  of $y = f(x - 3)$ sit? \teachesat{15--16}
  \answerto{Three to the right. Inside the bracket the number acts on the
  input, and the graph moves the opposite way to the sign.}
\item What does $y = 2f(x)$ do to the graph of $y = f(x)$? \teachesat{18--19}
  \answerto{Doubles every height, stretching it away from the horizontal axis.
  A point that was at $3$ is now at $6$; a point at $0$ does not move.}
\item Where does $y = (x + 4)^{2}$ have its lowest point? \teachesat{17--18}
  \answerto{At $x = -4$, where the bracket is zero.}
\item Is $2f(x) + 1$ the same as $2\bigl(f(x) + 1\bigr)$? \teachesat{22--23}
  \answerto{No. At $f(x) = 3$ the first is $7$ and the second is $8$. The moves
  are applied in the order they are written.}
\item What is $\sigma(0)$, where $\sigma(x) = \frac{1}{1 + e^{-x}}$?
  \teachesatone{25}
  \answerto{Exactly $\num{0.5}$, because $e^{0} = 1$ and $\frac{1}{1+1}$ is a
  half.}
\item Where does $\sigma(wx + b)$ cross a half? \teachesat{27--28}
  \answerto{At $x = -\frac{b}{w}$, which depends on the weight as well as on
  the bias.}
\item With $g(x) = 2x + 1$ and $f(x) = x^{2}$, evaluate $f(g(1))$ and
  $g(f(1))$. \teachesat{31--32}
  \answerto{$9$ and $3$. Composition is read from the inside out, and the two
  orders are different functions.}
\item Write the inverse of $f(x) = 2x + 1$. \teachesatone{36}
  \answerto{$f^{-1}(y) = \frac{y - 1}{2}$. Undo the operations in the reverse
  of the order they were done.}
\item Which of $x^{2}$, $e^{x}$, $-x$ and $\frac{x}{2}$ are strictly
  increasing? \teachesat{41--42}
  \answerto{$e^{x}$ and $\frac{x}{2}$. The first decreases for $x < 0$ and the
  third decreases everywhere.}
\item Four scores are $\num{2.0}$, $\num{1.5}$, $\num{0.5}$ and $-\num{1.0}$.
  Divide every one of them by $2$. Which is the largest now?
  \teachesat{44--45}
  \answerto{The first, as before. Halving every score cannot reorder them.}
\end{quiz}

% ============================================================
\section{One output, and only one}
\label{sec:F05-machine}
\index{function!domain of}
\index{function!one output per input}

\begin{fr}
A function is a machine. You put a number in at one end and exactly one number
comes out at the other. You have written thousands of them:

\begin{python}
def f(x):
    return x * x
\end{python}

Feed it $3$ and it hands back $9$. Feed it $3$ again tomorrow and it hands back
$9$ again \dash{} the machine has no memory and no mood, which is the whole of
what makes it a function rather than a process.

The paper notation is shorter and says the same thing: $f(x) = x^{2}$, read
\emph{f of x equals x squared}. The letter $x$ is a placeholder for whatever
you feed in, and $f(3)$ means \emph{the number this machine gives when you feed
it 3}.

What is $f(-3)$?
\nextframe
\end{fr}

\begin{fr}
\ans{$9$}
Squaring throws the sign away, so $-3$ and $3$ are handed the same answer. That
is allowed: a function may give the same output to two different inputs. What
it may not do is give two different outputs to the same input.

Here is a rule that breaks that condition. Take
\[ x^{2} + y^{2} = 18 \]
and read it as an instruction: \emph{given $x$, give me the $y$}.

Feed it $x = 3$. Is that a function of $x$?
\nextframe
\end{fr}

\begin{fr}
\ans{No}
At $x = 3$ the rule offers $y = 3$ and $y = -3$, because both satisfy
$9 + 9 = 18$. Two answers to one question is not a machine you can use: you
would have to ask which one it meant.

Drawn on paper the condition has a name you can apply by eye. Put a ruler
vertically anywhere on the graph. If it ever crosses the curve twice, the
picture is not the graph of a function. A circle fails at once; a parabola
never does, however far you slide the ruler.

Which of these three are functions of $x$?
\[ \text{(a)}\; y = x^{2} \qquad \text{(b)}\; y^{2} = x \qquad
   \text{(c)}\; y = \lvert x \rvert \]
\dotline
\nextframe
\end{fr}

\begin{fr}
\begin{ansblock}
(a) and (c).
\end{ansblock}

(b) fails: at $x = 4$ it offers $y = 2$ and $y = -2$, which is the circle's
fault under another name. Note that (c) gives $2$ for both $x = -2$ and
$x = 2$ and is perfectly respectable, for the reason frame $2$ gave.
\emph{Many-to-one} is a function; \emph{one-to-many} is not.

There is a second way a rule can fail, and it fails only at particular inputs
rather than everywhere. Consider
\[ f(x) = \frac{1}{x - 2} \]

What does it give at $x = 2$?
\nextframe
\end{fr}

\begin{fr}
\ans{Nothing at all}
There is no number to hand back, because the division has no answer. So $2$ is
not an input this machine accepts.

The set of inputs a function accepts is called its \textbf{domain}, and giving
it is part of giving the function. Three you will meet constantly:

\begin{center}
\begin{tabular}{ll}
\toprule
Function & Domain \\
\midrule
$\frac{1}{x - 2}$ & every $x$ except $2$ \\
$\sqrt{x}$ & $x \ge 0$, if the answer is to be a real number \\
$\ln x$ & $x > 0$, and Program~\ref{prog:F03} says why \\
\bottomrule
\end{tabular}
\end{center}

In code the domain is the set of arguments that do not raise, and the two
notions line up exactly: \code{math.sqrt(-1)} and \code{math.log(0)} both stop
the program, and they stop it because you asked a question outside the domain
rather than because anything went wrong.

Fill in the gap: the domain of $\sqrt{x - 5}$ is $x \ge \blank$.
\dotline
\nextframe
\end{fr}

\begin{fr}
\begin{ansblock}
$x \ge 5$.
\end{ansblock}

The bracket has to come out at zero or more, so $x - 5 \ge 0$. You have just
done section 3's work early, and you did it without being told a rule.

Now a distinction that costs working engineers real time, because the two
things look almost identical on the page. $f$ and $f(x)$ are not the same
object. $f$ is the machine; $f(x)$ is one number the machine produced. You can
pass a machine around, store it in a list, or hand it to something else that
will run it later. You cannot do any of that with a number.

In code the difference is two characters:

\begin{python}
act = relu
act = relu(x)
\end{python}

Say what each of those two lines leaves in \code{act}.
\nextframe
\end{fr}

\begin{fr}
\begin{ansblock}
The first leaves the function itself; the second leaves one number, computed
once, from whatever \code{x} happened to be at that moment.
\end{ansblock}

\begin{trapbox}
\textbf{The parentheses are the call, and writing them a line too early is one
of the standard ways to build a network that trains to nothing.} A layer that
was handed \code{relu(x)} rather than \code{relu} is holding a constant. It
does not fail, it does not warn, and it produces the same shaped output every
step \dash{} so the symptom arrives as a loss that will not move rather than as
an error that names the line.

The habit that prevents it is to read every name and ask which of the two you
are looking at. \emph{Machine, or number?} $f$, $\sigma$ and $\ln$ are
machines. $f(3)$, $\sigma(0)$ and $\ln 2$ are numbers.
\end{trapbox}
\end{fr}

\mermaidfig{f05-one-output}{Two rules and what each does with the input $3$.
The upper one returns a single value and is a function; the lower one returns
two and therefore answers no question you can ask it.}{the one-output
condition}

% ============================================================
\section{The graph, and what you can read off it}
\label{sec:F05-graph}
\index{graph!reading values off}
\index{function!increasing and decreasing}

\begin{fr}
The graph of $f$ is the set of all the points $(x, f(x))$, drawn. Nothing is
added by drawing it \dash{} the picture and the rule carry exactly the same
information \dash{} but the picture puts several questions within reach of the
eye that the rule only answers by arithmetic.

Build one. Take $f(x) = 2x + 1$ and fill in the outputs:

\begin{center}
\begin{tabular}{lccc}
\toprule
$x$ & $0$ & $1$ & $2$ \\
$f(x)$ & \blank & \blank & \blank \\
\bottomrule
\end{tabular}
\end{center}
\dotline
\nextframe
\end{fr}

\begin{fr}
\begin{ansblock}
$1$, $3$, $5$.
\end{ansblock}

Plot those three points and they lie on a straight line, climbing two units of
height for every one unit across. Join them and you have the graph.

Three things you can now read off without computing anything:

\begin{itemize}
\item \textbf{a value.} Go along to the $x$ you want, go up to the line, read
  across. That is $f(x)$.
\item \textbf{a crossing.} Where the line meets the horizontal axis, $f(x)$ is
  zero. That is where the equation $f(x) = 0$ is solved, and
  Program~\ref{prog:F06} does it by algebra.
\item \textbf{a direction.} Whether the curve is climbing or falling as you
  read left to right.
\end{itemize}

Where does $y = 2x + 1$ cross the horizontal axis?
\nextframe
\end{fr}

\begin{fr}
\ans{At $x = -\num{0.5}$}
Because $2x + 1 = 0$ needs $x = -\frac{1}{2}$. The line crosses once and never
again, which is a property of straight lines and not of graphs in general.

Now the direction. A function is \textbf{increasing} on an interval when moving
right along it never lowers the output, and \textbf{strictly increasing} when
moving right always raises it. Section 7 needs the strict version and needs it
exactly, so it is worth being fussy here: $2x + 1$ is strictly increasing
everywhere, because $2$ units of climb per unit across is $2$ units of climb
wherever you stand.

Over which values of $x$ is $y = x^{2}$ increasing?
\nextframe
\end{fr}

\begin{fr}
\ans{For $x \ge 0$}
And decreasing for $x \le 0$. At $x = -3$ the value is $9$; at $x = -1$ it is
$1$; at $x = 0$ it is $0$; and then it climbs again. The parabola falls into
its own bottom and comes back out.

\begin{warning}
\textbf{\enquote{Increasing} is a statement about an interval, never about a
function on its own.} A sentence that says \emph{$x^{2}$ is increasing} has
left out the only part that could make it true or false. Section 7 turns on
this distinction and will be no use to you if you carry it loosely.
\end{warning}

Here is one you will draw more often than any other. The rectified linear unit
is $\relu(x) = \max(0, x)$: it passes a positive number through untouched and
returns zero for anything negative.

Over which values of $x$ is its graph flat?
\nextframe
\end{fr}

\begin{fr}
\ans{For $x \le 0$}
To the left of the origin the graph lies along the horizontal axis; to the
right it is the line $y = x$, climbing one unit for every one across \dash{}
forty-five degrees if the axes are drawn to the same scale. The two halves meet
at the
origin in a corner rather than a curve.

That corner is worth a moment, because it is the entire reason the function is
there. A graph made of one straight piece is a straight line however long you
make it. A graph with a bend in it is not, and section 5 shows what that buys.

Compare $\relu(x)$ with $\lvert x \rvert$, which also has a corner at the
origin. How do the two differ?
\nextframe
\end{fr}

\begin{fr}
\ans{On the left half}
$\lvert x \rvert$ turns the negative side upwards and climbs away; $\relu$
flattens it to zero. On the right the two are identical.

These are the shapes the rest of the book leans on, and it is worth being able
to sketch each in three seconds:

\begin{center}
\begin{tabular}{ll}
\toprule
Rule & Shape \\
\midrule
$y = mx + c$ & a straight line, climbing at $m$ \\
$y = x^{2}$ & a parabola with its bottom at the origin \\
$y = \lvert x \rvert$ & a V with its point at the origin \\
$y = \relu(x)$ & the right half of that V, flat on the left \\
$y = \frac{1}{x}$ & two branches, neither touching either axis \\
\bottomrule
\end{tabular}
\end{center}

Program~\ref{prog:F07} adds the two that are missing from that list, the
exponential and the logistic, and it draws them rather than reciting them.
\end{fr}

% ============================================================
\section{Four moves you can do by eye}
\label{sec:F05-moves}
\index{graph!shifting}
\index{graph!stretching}
\index{graph!reflecting}

\begin{fr}
You now need one skill, and it is the one that turns a page of unfamiliar
formulas into five familiar shapes with numbers attached. Given the graph of
$f$, you should be able to draw the graph of anything built from it by adding,
multiplying and negating \dash{} without computing a single new value.

Start with the easy half. Take $f(x) = x^{2}$ and consider
\[ g(x) = f(x) + 2 \]

Every output is two bigger than it was, so every point on the graph is two
higher than it was, so the whole graph slides up by two. Check it at one point:
$f(3) = 9$, so $g(3) = 11$.

Where does the graph of $f(x) - 5$ sit relative to the graph of $f$?
\nextframe
\end{fr}

\begin{fr}
\ans{Five lower}
Same shape, slid down. A number added outside the bracket acts on the
\emph{output}, and it does exactly what it says.

Now the other half, and this one you should answer without working and without
checking. Write down the first thing that comes to you.

\textbf{The graph of $y = f(x)$ is in front of you. Which way does the graph of
$y = f(x - 3)$ sit, and by how much?}

Left or right, and a number.
\dotline
\nextframe
\end{fr}

\begin{fr}
\ans{Three to the \emph{right}}
Not left. Check it on the parabola, where you can see the answer rather than
argue about it. With $f(x) = x^{2}$ the lowest point sits at $x = 0$, because
that is where the bracket is zero. Now
\[ g(x) = f(x - 3) = (x - 3)^{2} \]
and \emph{its} bracket is zero at $x = 3$. So the bottom of the parabola \dash{}
the whole graph with it \dash{} has moved three to the right.

\begin{trapbox}
\textbf{Inside the bracket, the graph moves the opposite way to the sign, and
the reasoning that says otherwise is a true rule carried one place too far.}

Outside the bracket a minus really does move things down: $f(x) - 5$ sits five
lower, and you had that right one frame ago. What is different inside the
bracket is \emph{what the number is acting on}. It is not adjusting the answer;
it is adjusting the question. $g$ has to reach back three units to fetch the
value $f$ would have given, so to get the output $f$ gave at $0$ you must now
feed $g$ the input $3$.

Say it as a sentence and it stops being surprising: \emph{to see the old
picture you now have to stand three further along.}

This reasoning has arrived under other signs already. Program~\ref{prog:F01}
has $\sqrt{a+b} \ne \sqrt{a} + \sqrt{b}$; Program~\ref{prog:F02} has
$(a+b)^{2} \ne a^{2} + b^{2}$; Program~\ref{prog:F03} has no rule for
$\ln(a+b)$; Program~\ref{prog:F04} has
$\sum a_i b_i \ne (\sum a_i)(\sum b_i)$. Each time the fix is the same
question: \emph{what is this operation, and what is it being applied to?}
\end{trapbox}
\end{fr}

\begin{fr}
Confirm it with numbers rather than with the argument, because the argument is
the thing you just got wrong and a table is not going to lie to you.

\begin{center}
\begin{tabular}{lccccc}
\toprule
$x$ & $1$ & $2$ & $3$ & $4$ & $5$ \\
$(x - 3)^{2}$ & $4$ & $1$ & $0$ & $1$ & $4$ \\
\bottomrule
\end{tabular}
\end{center}

The low point is at $x = 3$ and the shape either side of it is the parabola you
started with. Three to the right, confirmed.

Where is the lowest point of $y = (x + 4)^{2}$?
\nextframe
\end{fr}

\begin{fr}
\ans{At $x = -4$}
The bracket is zero there, and the sign is a plus, so the graph went the other
way again \dash{} four to the \emph{left}. State it once and keep it:

\begin{center}
\begin{tabular}{ll}
\toprule
Written & Does \\
\midrule
$f(x) + c$ & slides the graph up by $c$ \\
$f(x - c)$ & slides the graph right by $c$ \\
\bottomrule
\end{tabular}
\end{center}

Two moves left, and they are the multiplying versions of the same pair. What
does the graph of $y = 2f(x)$ look like, compared with the graph of $y = f(x)$?
\nextframe
\end{fr}

\begin{fr}
\ans{Twice as tall}
Every output is doubled, so every point is twice as far from the horizontal
axis as it was. A point that was at height $3$ is now at $6$; a point that was
at height $-1$ is now at $-2$; and a point that was at $0$ has not moved at
all, because twice nothing is nothing. The graph is stretched away from the
horizontal axis, and the axis itself is pinned.

Now the one inside the bracket. What does $y = f(2x)$ do?
\nextframe
\end{fr}

\begin{fr}
\ans{Squashes it towards the vertical axis, by two}
Backwards again, and for the same reason as before. To get the value $f$ gave
at $6$ you now feed $3$, because the machine doubles what you hand it before
$f$ ever sees it. Everything therefore happens at half the distance out, and
the graph is compressed horizontally rather than stretched.

That is a genuinely useful sanity check when you meet an unfamiliar formula:
\emph{a big number inside the bracket makes the picture narrower, not wider.}

Two reflections and the section is complete. What do $y = -f(x)$ and
$y = f(-x)$ each do?
\nextframe
\end{fr}

\begin{fr}
\begin{ansblock}
$-f(x)$ flips the graph top to bottom; $f(-x)$ flips it left to right.
\end{ansblock}

Both are the $a = -1$ case of the two stretches, and both obey the same rule
about where the number is written. Outside the bracket the \emph{output}
changes sign, so up becomes down. Inside the bracket the \emph{input} changes
sign, so the value that used to be at $x$ is now at $-x$.

Take $f(x) = x^{2}$. What does $f(-x)$ look like?
\nextframe
\end{fr}

\begin{fr}
\ans{Exactly the same}
Because $(-x)^{2} = x^{2}$. The parabola is its own mirror image in the
vertical axis, so the move happens and nothing visible results. A function with
that property is called \textbf{even}; $\lvert x \rvert$ is another, and
$\cos$, which Program~\ref{prog:F08} draws, is a third.

One thing left, and it is where the marks go in an examination and where the
bugs go in a model. \textbf{The moves do not commute.} Is $2f(x) + 1$ the same
function as $2\bigl(f(x) + 1\bigr)$?
\nextframe
\end{fr}

\begin{fr}
\ans{No}
Take a point where $f(x) = 3$. The first gives $2 \times 3 + 1 = 7$; the second
gives $2 \times (3 + 1) = 8$. The first stretches and then slides; the second
slides and then stretches, and the slide gets stretched along with everything
else.

So read the expression the way arithmetic reads it \dash{} brackets first,
then multiply, then add \dash{} and apply the moves in that order. There is no
extra rule to learn here; the ordinary order of operations is already the
answer.
\end{fr}

\mermaidfig{f05-four-moves}{Where a number is written decides which way it
acts. Outside the bracket it adjusts the answer; inside it adjusts the
question, and the graph moves the other way.}{the four graph moves}

% ============================================================
\section{Weight and bias are a stretch and a shift}
\label{sec:F05-weight-bias}
\index{function!logistic}
\index{bias!as a shift}

\begin{fr}
Everything in the previous section was practice for this. Here is the most
common expression in machine learning:
\[ y = wx + b \]

Read it as two of the four moves applied to the plainest function there is,
$f(x) = x$. The weight $w$ multiplies the output, so it stretches the line away
from the horizontal axis: a $w$ bigger than $1$ makes it steeper, a $w$ between
$0$ and $1$ flattens it, and a negative $w$ turns it over as well. The bias $b$
is added outside,
so it slides the whole line up. That is all a linear unit is: a stretch and a
shift, with names from a different vocabulary.

At $w = 3$ and $b = -1$, what is $y$ when $x = 2$?
\nextframe
\end{fr}

\begin{fr}
\ans{$5$}
From $3 \times 2 - 1$. Now put the same two moves inside a curve, because that
is where the vocabulary starts to earn its keep. The \textbf{logistic
function} is
\[ \sigma(x) = \frac{1}{1 + e^{-x}} \]
and Program~\ref{prog:F07} derives its shape; here you only need to be able to
read it.

One value you should know without computing: $\sigma(0) = \num{0.5}$, exactly,
because $e^{0} = 1$ and $\frac{1}{1+1}$ is a half. Around that centre it runs
between $0$ and $1$ and reaches neither:

\begin{center}
\begin{tabular}{rc}
\toprule
$x$ & $\sigma(x)$ \\
\midrule
$-4$ & \val{f05.sig.m4} \\
$-2$ & \val{f05.sig.m2} \\
$-1$ & \val{f05.sig.m1} \\
$0$ & \val{f05.sig.0} \\
$1$ & \val{f05.sig.1} \\
$2$ & \val{f05.sig.2} \\
$4$ & \val{f05.sig.4} \\
\bottomrule
\end{tabular}
\end{center}

Look at the table rather than at the formula. What is the relationship between
$\sigma(-2)$ and $\sigma(2)$?
\nextframe
\end{fr}

\begin{fr}
\begin{ansblock}
They add to $1$: $\val{f05.sig.m2} + \val{f05.sig.2} = 1$.
\end{ansblock}

So $\sigma(-x) = 1 - \sigma(x)$, and the curve has a half-turn symmetry about
the point $\bigl(0, \frac{1}{2}\bigr)$. Every row of that table comes in such a
pair, which halves what there is to remember.

Now apply the two moves to the \emph{input} instead, which is what a layer
actually does:
\[ \sigma(wx + b) \]

The weight is inside the bracket, so by the previous section it squashes the
curve horizontally. What does that do to how the curve looks?
\nextframe
\end{fr}

\begin{fr}
\ans{It makes it steeper}
A bigger weight compresses the same rise into a narrower stretch of $x$, so the
curve turns from a gentle slope into something approaching a step.

That is worth a number rather than an adjective. Measure the width of the band
in which the curve does its work \dash{} from $\sigma = \num{0.1}$ up to
$\sigma = \num{0.9}$:

\begin{center}
\begin{tabular}{lcc}
\toprule
Weight & Width of the band & \\
\midrule
$w = \num{0.5}$ & \val{f05.band.whalf} & twice as wide \\
$w = 1$ & \val{f05.band.w1} & \\
$w = 5$ & \val{f05.band.w5} & a fifth as wide \\
\bottomrule
\end{tabular}
\end{center}

The band is divided by exactly $w$, which is the horizontal squash of frame
$20$ with a measurement attached. At $w = 5$ the curve does the whole of its
work inside a window of $x$ less than one unit wide.

Where does $\sigma(wx + b)$ cross a half?
\nextframe
\end{fr}

\begin{fr}
\ans{At $x = -\frac{b}{w}$}
Because the crossing happens where the bracket is zero, and $wx + b = 0$ gives
$x = -\frac{b}{w}$. Not at $x = b$, and not at $x = -b$.

\begin{trapbox}
\textbf{The bias is not a position.} It is tempting to read $b$ as
\emph{where this unit fires}, because it is the number that moves the curve
sideways and it is the number you can see in the configuration. But the
distance it moves the curve is $\frac{b}{w}$, so the same bias means different
things at different weights.

At $b = -2$ the curve crosses a half at $x = 2$ when $w = 1$, and at
$x = \num{0.4}$ when $w = 5$. One number, unchanged; the threshold it sets,
five times nearer the origin.

This is the inside-the-bracket rule of section 3 with a second complication
stacked on it: the number acts on the input, so it moves the graph the other
way, \emph{and} it is scaled by whatever else is acting on the input. Read
$-\frac{b}{w}$, never $b$.
\end{trapbox}
\end{fr}

\begin{fr}
One question is still open. If the weight sets the steepness and the crossing
point is $-\frac{b}{w}$, why carry a bias at all \dash{} why not set $b = 0$
and let the weight do the work?

Here is the smallest answer. Take three measurements:
\[ (1, 3) \qquad (2, 5) \qquad (3, 7) \]
The line $y = 2x + 1$ passes through all three exactly: $2 + 1$, $4 + 1$,
$6 + 1$.

Now forbid the bias, so the model is $y = wx$ and nothing else. The line of
that form that makes the sum of the squared misses smallest \dash{} the
\emph{least-squares} line, which Program~\ref{prog:P08} derives \dash{} has
slope $\val{f05.nobias.slope}$, and its worst miss against the three points is
$\val{f05.nobias.worst}$.

Why can no line of the form $y = wx$ fit these three points, whatever $w$ you
choose?
\nextframe
\end{fr}

\begin{fr}
\begin{ansblock}
Because every such line passes through the origin, and these points do not lie
on any line that does.
\end{ansblock}

$y = wx$ is worth $0$ at $x = 0$, always, for every weight. The data says the
value at $x = 0$ should be $1$. No amount of steepness reaches a point the
family cannot express, so the bias is not a refinement \dash{} it is the
difference between a family of lines through one fixed point and a family of
all lines.

\begin{aibox}
Layer normalisation ends with $\gamma \hat{x} + \beta$: it standardises the
values, and then applies a learned stretch and a learned shift to put them
back where the network wants them. Those are exactly the two outside-the-bracket
moves of section 3, one of each, and they are parameters the model trains.

So the normalisation is not the end of that layer's story. Having removed the
scale and the offset, the layer immediately learns a scale and an offset to
reapply \dash{} which sounds circular until you notice that the ones it learns
do not depend on the batch it happens to be looking at, and the ones it removed
did.
\end{aibox}
\end{fr}

% ============================================================
\section{Machines wired in series}
\label{sec:F05-composition}
\index{function!composition of}

\begin{fr}
Feed the output of one machine straight into the next and you have a
\textbf{composition}, written $f(g(x))$ and read \emph{f of g of x}.

There is no new idea here, only a reading order: work from the inside out, the
way you would run the code. With $g(x) = 2x + 1$ and $f(x) = x^{2}$:
\[ f(g(1)) = f(2 \times 1 + 1) = f(3) = 9 \]

Now evaluate $g(f(1))$.
\nextframe
\end{fr}

\begin{fr}
\ans{$3$}
Because $f(1) = 1$ and then $g(1) = 3$. Nine against three from the same two
machines and the same input: \textbf{the order is part of the function}, and
$f(g(x))$ and $g(f(x))$ are two different functions that happen to be built
from the same parts.

Here is what that buys, and it is the reason a network has activation functions
in it at all. Wire two linear units in series:
\[ y = w_2(w_1 x + b_1) + b_2 = (w_2 w_1)x + (w_2 b_1 + b_2) \]

Look at what came out. Multiply out the bracket and the whole thing is
$Wx + B$ for a single weight $W = w_2 w_1$ and a single bias
$B = w_2 b_1 + b_2$.

What does that say about a stack of linear layers with nothing between them?
\nextframe
\end{fr}

\begin{fr}
\begin{ansblock}
It is a single linear layer, however many you stack.
\end{ansblock}

Two hundred of them collapse to one, with one weight and one bias, and the
extra $199$ have bought nothing but arithmetic. Depth on its own is not a
capability.

\begin{aibox}
That is the whole argument for putting a $\relu$ between the layers. Its graph
has a corner in it, from frame $12$, and no amount of stretching and sliding
produces a corner from a straight line \dash{} so
$\relu(w_2\,\relu(w_1 x + b_1) + b_2)$ does not collapse, and cannot be
rewritten as $Wx + B$.

The activation function is not there to squash values into a range or to be
biologically suggestive. It is there because without it the composition is
provably a waste of parameters, and the proof is the two lines of algebra you
have just read.
\end{aibox}

A network of $32$ layers is $32$ of these machines nested. How would you
differentiate that?
\nextframe
\end{fr}

\begin{fr}
\begin{ansblock}
With the chain rule, which is what Program~\ref{prog:F12} is for.
\end{ansblock}

That is the hinge of this book, and this frame is where the two halves meet.
Backpropagation is not an algorithm invented for neural networks; it is the
chain rule applied to a composition, with the intermediate results kept rather
than recomputed. You have just built the composition. Programs~\ref{prog:F11}
and~\ref{prog:F12} supply the differentiation.

One caution before leaving. You may not always wire two machines together. What
has to be true of $g$ before $f(g(x))$ makes sense?
\nextframe
\end{fr}

\begin{fr}
\begin{ansblock}
Every output of $g$ has to be an input $f$ accepts \dash{} $g$'s outputs must
lie inside $f$'s domain.
\end{ansblock}

The two compositions of $\ln$ and $\sigma$ show both cases. $\ln(\sigma(x))$ is
fine for every $x$, because $\sigma$ produces strictly positive numbers and
that is exactly what $\ln$ demands. Turn them round and $\sigma(\ln x)$ needs
$x > 0$, because $\ln$ is the one with the restricted domain and it is now
first in the queue.

\begin{note}
$\ln(\sigma(x))$ is worth recognising when you meet it. It is the log-sigmoid,
and it appears in a binary cross-entropy whenever the input is a logit.
Composing the two is safe mathematically, as above, and unsafe in floating
point for large negative $x$,
where $\sigma$ underflows to zero before $\ln$ ever sees it.
Program~\ref{prog:P02} shows why the two operations are done as one.
\end{note}
\end{fr}

% ============================================================
\section{Undoing a function}
\label{sec:F05-inverse}
\index{function!inverse of}

\begin{fr}
If a function is a machine, an inverse is the machine run backwards: hand it an
output and it tells you which input produced it. Written $f^{-1}$.

Finding one is Program~\ref{prog:F06}'s work done informally: undo the
operations in the reverse of the order they were applied. For $f(x) = 2x + 1$
the machine doubles and then adds one, so the inverse subtracts one and then
halves:
\[ f^{-1}(y) = \frac{y - 1}{2} \]
Check it: $f(3) = 7$, and $f^{-1}(7) = \frac{6}{2} = 3$.

What is the inverse of $f(x) = x^{3}$?
\nextframe
\end{fr}

\begin{fr}
\ans{The cube root}
$f^{-1}(y) = \sqrt[3]{y}$, and it works for negative numbers as well, because
cubing keeps the sign.

Squaring does not, and that is the whole of the condition. An inverse exists
only when the function is \textbf{one-to-one}: different inputs must give
different outputs, or there is no way to say which input an output came from.
$f(x) = x^{2}$ fails, because $9$ could have come from $3$ or from $-3$ \dash{}
which is frame $2$'s observation arriving as an obstruction. Restrict the
domain to $x \ge 0$ and the inverse exists again, and it is $\sqrt{y}$.

If you have the graph of $f$ in front of you, how do you draw the graph of
$f^{-1}$?
\nextframe
\end{fr}

\begin{fr}
\begin{ansblock}
Reflect it in the line $y = x$.
\end{ansblock}

Because reflecting in that line swaps the two coordinates of every point, and
swapping the coordinates is precisely what an inverse does: the point
$(3, 7)$ on the graph of $f$ becomes the point $(7, 3)$ on the graph of
$f^{-1}$.

You have already used a pair like this. Program~\ref{prog:F03} spent a whole
program on $\ln$ and $e^{x}$, which are inverses of each other:
$\ln(e^{x}) = x$ and $e^{\ln x} = x$ for $x > 0$. Their graphs are mirror
images in $y = x$, and the restricted domain of one is the restricted range of
the other.

The logistic has an inverse too. Solve $p = \sigma(x)$ for $x$.
\nextframe
\end{fr}

\begin{fr}
\begin{ansblock}
\[ x = \ln\!\left(\frac{p}{1 - p}\right) \]
\end{ansblock}

It has a name, the \textbf{logit}, and $\frac{p}{1-p}$ has one too \dash{} the
\emph{odds}. So a logit is the logarithm of the odds, which is the sentence
that makes the word mean something rather than being a label for whatever comes
out of the last layer.

At $p = \num{0.9}$ the odds are nine to one and the logit is
$\val{f05.logit.p9}$. Push that back through $\sigma$ and you get
$\num{0.9}$ again, to within $\val{f05.roundtrip.err}$ \dash{} which is the
round trip working as well as binary64 allows and not one bit better.

Last question of the section, and it is about notation rather than about
functions. Is $f^{-1}(x)$ the same thing as $\frac{1}{f(x)}$?
\nextframe
\end{fr}

\begin{fr}
\ans{No}
They are unrelated. With $f(x) = 2x + 1$, the inverse at $7$ is $3$ and the
reciprocal at $7$ is $\frac{1}{15}$.

\begin{trapbox}
\textbf{The $-1$ in $f^{-1}$ is not an exponent}, even though it is written
exactly where an exponent goes and even though every other superscript in the
book is one. $f^{2}(x)$ is ambiguous in the wild \dash{} some authors mean
$f(f(x))$ and some mean $(f(x))^{2}$ \dash{} and the trigonometric functions
make the collision worse rather than better, because $\sin^{2}x$ conventionally
means $(\sin x)^{2}$ while $\sin^{-1}x$ means the inverse.

This book writes $\arcsin$ for the inverse, and Program~\ref{prog:F08} says so
again where it matters. Where you are reading somebody else's paper, the only
safe move is to check what they do with a case you can evaluate.
\end{trapbox}
\end{fr}

% ============================================================
\section{The order a function keeps, and the one it does not}
\label{sec:F05-monotone}
\index{function!strictly increasing}
\index{softmax!temperature and}

\begin{fr}
One property is left, and the rest of this section is what it buys.

A function is \textbf{strictly increasing} when moving right along it always
raises the output: whenever $a < b$, it is also true that $f(a) < f(b)$. Not
\emph{sometimes}, not \emph{on average}, and not \emph{eventually} \dash{}
every pair, everywhere it is defined.

Adding five is strictly increasing. Doubling is. Halving is, and so is dividing
by any positive number, because dividing by a positive number keeps every
comparison the way it was.

Which of these four are strictly increasing over all the values they accept?
\[ x^{2} \qquad e^{x} \qquad -x \qquad \frac{x}{2} \]
\dotline
\nextframe
\end{fr}

\begin{fr}
\begin{ansblock}
$e^{x}$ and $\frac{x}{2}$.
\end{ansblock}

$x^{2}$ fails on the left half, where $-3 < -1$ but $9 > 1$; the warning at
frame $11$ was about exactly this. And $-x$ fails everywhere and fails
systematically \dash{} it is strictly \emph{decreasing}, which is a property in
its own right and not merely the absence of one.

Now answer this without working it out, and write the first thing that comes
to you.

\textbf{Your sampler is producing text that repeats itself, so you raise the
temperature from $\num{1.0}$ to $\num{2.0}$. Can that change which token the
model considers most likely?}

Yes or no.
\dotline
\nextframe
\end{fr}

\begin{fr}
\ans{No. It cannot, at any temperature}
Here is the evidence before the argument. Take four scores \dash{} logits, in
the vocabulary of the last layer \dash{} and turn them into probabilities at
three temperatures with
\[ \softmax(z)_j = \frac{e^{z_j / T}}{\sum_k e^{z_k / T}} \]

\begin{center}
\begin{tabular}{rccc}
\toprule
Score $z$ & $T = \num{0.5}$ & $T = \num{1.0}$ & $T = \num{2.0}$ \\
\midrule
$\num{2.0}$ & \val{f05.sm.t05.p1} & \val{f05.sm.t10.p1} & \val{f05.sm.t20.p1} \\
$\num{1.5}$ & \val{f05.sm.t05.p2} & \val{f05.sm.t10.p2} & \val{f05.sm.t20.p2} \\
$\num{0.5}$ & \val{f05.sm.t05.p3} & \val{f05.sm.t10.p3} & \val{f05.sm.t20.p3} \\
$-\num{1.0}$ & \val{f05.sm.t05.p4} & \val{f05.sm.t10.p4} & \val{f05.sm.t20.p4} \\
\bottomrule
\end{tabular}
\end{center}

Read down each column: whatever the temperature, the largest entry is the
one in the top row. Now read along each row: every single number moved.

\begin{trapbox}
\textbf{\enquote{Raising the temperature made it pick a different top token} is
impossible before sampling and expected after it}, and keeping those two apart
is the whole of the fix.

What temperature changes is enormous, and it is easy to look at the wrong end
of the distribution and conclude that it changes little. Across that table the
most likely token falls from \val{f05.sm.t05.top} to \val{f05.sm.t20.top}, a
factor of \val{f05.sm.top.ratio}. The \emph{least} likely one rises from
\val{f05.sm.t05.p4} to \val{f05.sm.t20.p4}, more than
\val{f05.sm.tail.floor}-fold. Temperature is a knob on the tail, which is
exactly what \enquote{more creative} means.

So the observable consequence is real: the chance of drawing something other
than the top token goes from \val{f05.sm.t05.other} to
\val{f05.sm.t20.other}. What has not changed, and cannot, is which token is at
the top. If you are reading the argmax you will see the same answer at every
temperature you try, and if you conclude from that that the setting is being
ignored you will go looking for a bug that is not there.
\end{trapbox}
\end{fr}

\begin{fr}
Now the reason, which is one sentence and a corollary.

\textbf{A strictly increasing function keeps the order of a list.} If
$z_1 > z_2$ then $f(z_1) > f(z_2)$, by the definition and nothing else; apply
that to every pair and the whole ranking survives. In particular the entry that
was largest is still largest, so the position of the maximum \dash{} the
$\argmax$ \dash{} does not move.

That is the entire result. The rest is checking that the pieces qualify.

Which strictly increasing functions does the softmax apply to those scores,
before it divides?
\nextframe
\end{fr}

\begin{fr}
\begin{ansblock}
Two: dividing by $T$, and then exponentiating.
\end{ansblock}

Dividing by $T$ is strictly increasing for every $T > 0$, and $e^{x}$ is
strictly increasing everywhere, which you settled at frame $42$. Neither can
disturb the ranking, so the ranking that arrives at the division is the ranking
the scores had.

And the division itself cannot disturb it either: every entry is divided by the
same positive total, which is a stretch of the whole list towards zero and
leaves every comparison the way it was.

So $\argmax \softmax(z / T) = \argmax z$, for every $T > 0$ and every $z$. That
is exact in the algebra, and no vector of scores breaks it. Push $T$ high
enough and the arithmetic flattens every entry to the same float, at which
point the comparison has nothing left to compare \dash{} which is a limit of
binary64 and not of the result.

Program~\ref{prog:F03} left you a habit of working with logarithms of
probabilities rather than probabilities. Does taking a logarithm move the
$\argmax$?
\nextframe
\end{fr}

\begin{fr}
\ans{No}
$\ln$ is strictly increasing on its whole domain, and probabilities are
positive, so the largest probability has the largest logarithm.

That is why a decoder never has to leave log space to choose a token, and it is
worth seeing as the payoff it is: Program~\ref{prog:F03} adopted logarithms
because products underflow, and this frame says the adoption costs nothing at
the point where a decision is made. You compare log-probabilities, you pick the
largest, and it is the same token you would have picked from the
probabilities \dash{} without ever computing them.

One question left, and it is the one that decides whether you have the
property or only the example. What does a strictly \emph{decreasing} function
do to the $\argmax$?
\nextframe
\end{fr}

\begin{fr}
\begin{ansblock}
It moves it to where the $\argmin$ was.
\end{ansblock}

Every comparison is reversed, so the largest becomes the smallest and the
smallest becomes the largest. The order is not destroyed \dash{} it is turned
round, which is a different thing and is just as useful once you know which you
have.

\begin{aibox}
Negating is the simplest strictly decreasing function there is, and it is the
reason two sentences you have read a thousand times are one sentence.
\emph{Maximise the likelihood} and \emph{minimise the negative log-likelihood}
name the same parameters, because $\ln$ is strictly increasing and negation
strictly decreasing, so the composition of the two turns the $\argmax$ of the
likelihood into the $\argmin$ of the loss and moves it nowhere else.

That is also why the loss you train on may be scaled by a \emph{positive}
constant, shifted by a constant, or \dash{} if it is positive \dash{} have its
logarithm taken, and the best model is the same model. It is not why
the \emph{gradients} are the same, which is a separate question and belongs to
Program~\ref{prog:F12}.
\end{aibox}
\end{fr}

\mermaidfig{f05-order-kept}{Two strictly increasing steps, applied one after
the other to the same four scores. Neither can reorder them, so the largest
survives as the largest.}{order kept by increasing maps}

% ============================================================
\begin{summarybox}
\sumitem{1--4}{\result{\textbf{A function is a machine with one output per input.} Two
  outputs for one input is not a function \dash{} $x^{2} + y^{2} = 18$ offers
  $3$ and $-3$ at $x = 3$ \dash{} and the vertical ruler is that condition
  applied by eye. Two inputs sharing one output is perfectly ordinary}.}
\sumitem{4--5}{\result{The \emph{domain} is the set of inputs the machine accepts, and
  giving it is part of giving the function. It is the same set as the
  arguments that do not raise: $\frac{1}{x-2}$ excludes $2$, $\sqrt{x}$ needs
  $x \ge 0$, $\ln x$ needs $x > 0$}.}
\sumitem{6--7}{\result{$f$ is the machine and $f(x)$ is one number it produced.
  \code{act = relu} and \code{act = relu(x)} differ by two characters and by
  everything else, and the second fails as a loss that will not move rather
  than as an error}.}
\sumitem{8--10}{\result{The graph is the set of points $(x, f(x))$ drawn, and a value,
  a crossing and a direction are all readable off it without arithmetic}.}
\sumitem{10--11}{\result{\textbf{\enquote{Increasing} is a statement about an
  interval.} $x^{2}$ increases for $x \ge 0$ and decreases for $x \le 0$, so a
  sentence about $x^{2}$ alone has left out the part that decides it}.}
\sumitem{12--13}{\result{The five shapes the book leans on: line, parabola, V,
  $\relu$'s half-V and $\frac{1}{x}$. $\relu$ is flat for $x \le 0$ and meets
  the other half in a corner}.}
\sumitem{14--16}{\result{\textbf{Where the number is written decides which way it
  acts.} $f(x) + c$ slides the graph up by $c$, and $f(x - c)$ slides it
  \emph{right} by $c$. Outside the bracket it adjusts the answer; inside it
  adjusts the question, so you must stand $c$ further along to see the old
  picture}.}
\sumitem{18--21}{\result{$a f(x)$ stretches away from the horizontal axis by $a$;
  $f(ax)$ squashes towards the vertical axis by $a$; $-f(x)$ flips top to
  bottom and $f(-x)$ flips left to right. The inside pair go backwards, for
  the same reason the shift does}.}
\sumitem{22--23}{\result{The moves do not commute: $2f(x) + 1$ and
  $2\bigl(f(x)+1\bigr)$ are $7$ and $8$ where $f(x) = 3$. Ordinary precedence
  is the whole rule}.}
\sumitem{24--26}{\result{$y = wx + b$ is a stretch and a shift applied to $f(x) = x$.
  $\sigma(0) = \num{0.5}$ exactly, and $\sigma(-x) = 1 - \sigma(x)$, so the
  curve has a half-turn symmetry about $\bigl(0, \frac{1}{2}\bigr)$}.}
\sumitem{27--28}{\result{In $\sigma(wx + b)$ the weight sets the steepness: the band
  from $\num{0.1}$ to $\num{0.9}$ is \val{f05.band.w1} wide at $w = 1$ and
  \val{f05.band.w5} at $w = 5$, divided by exactly $w$.
  \textbf{The bias is not a position} \dash{} the crossing is at
  $-\frac{b}{w}$, so one bias sets five different thresholds at five different
  weights}.}
\sumitem{29--30}{\result{Without a bias every line passes through the origin, and no
  choice of weight reaches a point that family cannot express: the best
  bias-free fit to $(1,3)$, $(2,5)$, $(3,7)$ has slope
  \val{f05.nobias.slope} and misses by \val{f05.nobias.worst}, where
  $y = 2x + 1$ is exact. Layer normalisation's $\gamma$ and $\beta$ are the
  same two moves, learned}.}
\sumitem{31--33}{\result{$f(g(x))$ is read from the inside out, and the order is part
  of the function. \textbf{Two linear layers in series are one linear layer},
  $w_2(w_1x + b_1) + b_2 = (w_2w_1)x + (w_2b_1 + b_2)$, which is the whole
  argument for putting a function with a corner in it between them}.}
\sumitem{34--35}{\result{A deep network is a composition and nothing else, which is
  why Program~\ref{prog:F12}'s chain rule is this book's hinge. Two machines
  compose only when the first's outputs lie in the second's domain}.}
\sumitem{36--38}{\result{An inverse runs the machine backwards and exists only for a
  one-to-one function; $x^{2}$ has none over all the reals, for frame $2$'s
  reason. Its graph is the original reflected in $y = x$}.}
\sumitem{39--40}{\result{The logit $\ln\frac{p}{1-p}$ inverts the logistic \dash{} it
  is the logarithm of the odds, $\val{f05.logit.p9}$ at $p = \num{0.9}$, and
  the round trip closes to $\val{f05.roundtrip.err}$. \textbf{The $-1$ in
  $f^{-1}$ is not an exponent.}}}
\sumitem{41--45}{\result{\textbf{A strictly increasing function keeps the order of a
  list, and therefore cannot move the $\argmax$.} Dividing by a positive $T$
  qualifies, $e^{x}$ qualifies, and dividing every entry by one positive total
  qualifies \dash{} so $\argmax \softmax(z/T) = \argmax z$ for every $T > 0$,
  exactly and not approximately}.}
\sumitem{43}{\result{Temperature changes the distribution a great deal and the top of
  it least. Over $T$ from $\num{0.5}$ to $\num{2.0}$ the most likely token
  falls by a factor of \val{f05.sm.top.ratio} and the least likely rises more
  than \val{f05.sm.tail.floor}-fold; the chance of drawing something other
  than the top token goes from \val{f05.sm.t05.other} to
  \val{f05.sm.t20.other}. What never moves is which token is at the top}.}
\sumitem{46--47}{\result{$\ln$ is strictly increasing, so a decoder may compare
  log-probabilities and never compute a probability. A strictly
  \emph{decreasing} function swaps $\argmax$ and $\argmin$, which is why
  maximising a likelihood and minimising its negative logarithm are one
  instruction}.}
\end{summarybox}

\canyou

% ============================================================
\begin{testexercises}
\item Is $y^{3} = x$ a function of $x$? Is $y^{4} = x$?
  \answerto{The first is: every $x$ has exactly one real cube root. The second
  is not: at $x = 16$ it offers $y = 2$ and $y = -2$.}
\item Give the domain of $f(x) = \frac{\sqrt{x - 1}}{x - 4}$.
  \answerto{$x \ge 1$ and $x \ne 4$. The root needs a non-negative bracket and
  the division needs a non-zero one, and both conditions apply.}
\item A colleague builds a layer with \code{activation=sigmoid(x)} and the
  model trains to a flat loss. What has gone wrong?
  \answerto{The parentheses. The layer was handed one number rather than the
  function, so nothing it does depends on its own input.}
\item Where does $y = 3x - 6$ cross each axis?
  \answerto{At $x = 2$ on the horizontal axis and at $y = -6$ on the vertical
  one. Set the other coordinate to zero in each case.}
\item Over which values of $x$ is $y = -x^{2}$ increasing?
  \answerto{For $x \le 0$. Negating the parabola turns it upside down and
  reverses both halves.}
\item The graph of $y = f(x)$ is given. Describe the graphs of $y = f(x) - 4$
  and $y = f(x + 4)$.
  \answerto{Four lower, and four to the \emph{left}. The second moves the
  opposite way to its sign because the number is inside the bracket.}
\item Where is the lowest point of $y = (x - 7)^{2} + 2$?
  \answerto{At $(7, 2)$. The bracket is zero at $x = 7$ and the outside $+2$
  lifts the whole curve.}
\item Describe the graph of $y = \frac{1}{2}f(3x)$ in terms of the graph of
  $y = f(x)$.
  \answerto{Half as tall and a third as wide. The $\frac{1}{2}$ is outside and
  halves every height; the $3$ is inside and compresses horizontally by three.}
\item Where does $\sigma(2x - 6)$ cross a half, and is it steeper or gentler
  than $\sigma(x)$?
  \answerto{At $x = 3$, from $-\frac{-6}{2}$, and twice as steep. Reading the
  bias as the crossing point would have given $-6$.}
\item With $g(x) = x + 1$ and $f(x) = 3x$, evaluate $f(g(2))$ and $g(f(2))$.
  \answerto{$9$ and $7$. The first adds then triples; the second triples then
  adds.}
\item Three linear layers are stacked with nothing between them. How many
  parameters does the resulting function actually have?
  \answerto{Two \dash{} one weight and one bias \dash{} because the
  composition multiplies out to $Wx + B$ however many layers went into it.}
\item Write the inverse of $f(x) = \frac{x}{3} - 2$.
  \answerto{$f^{-1}(y) = 3(y + 2)$. Undo the subtraction first and the
  division second, which is the reverse of the order they were applied.}
\item A probability is $\num{0.25}$. What are its odds, and what is its logit?
  \answerto{Odds of $1$ to $3$, and a logit of $\ln \frac{1}{3}$, which is
  negative because the probability is below a half.}
\item Four scores are $\num{3.0}$, $\num{1.0}$, $\num{0.0}$ and
  $-\num{2.0}$. Which is largest after you add $10$ to every one of them, and
  which after you exponentiate every one?
  \answerto{The first, both times. Adding a constant and exponentiating are
  both strictly increasing, so neither can reorder the list.}
\item Your sampler runs at $T = \num{2.0}$ and you want the most likely token
  instead. Does lowering $T$ get it for you?
  \answerto{Not reliably, and not because of $T$. The most likely token is the
  same at every temperature; lowering $T$ only makes it more likely to be the
  one drawn. Taking the $\argmax$ gets it with certainty at any $T$.}
\end{testexercises}

\begin{furtherproblems}
\item Give a rule that assigns exactly one output to every input and is
  nonetheless useless as a graph to read values off.
  \answerto{For instance \emph{$f(x) = 1$ if $x$ is rational and $0$
  otherwise}. It is a perfectly good function and its graph cannot be drawn,
  which is a reminder that the picture is a convenience and not the
  definition.}
\item $f$ has domain $x \ge 0$. What is the domain of $f(x - 3)$?
  \answerto{$x \ge 3$. The domain shifts with the graph, and in the same
  direction.}
\item Show that $f(x) = \lvert x \rvert$ is even and that $f(x) = x^{3}$ is
  odd, where \emph{odd} means $f(-x) = -f(x)$.
  \answerto{$\lvert -x \rvert = \lvert x \rvert$, so the graph is unchanged by
  a left-right flip. $(-x)^{3} = -x^{3}$, so a left-right flip and a
  top-to-bottom flip give the same picture.}
\item By how much, and in which direction, does the graph of
  $y = f(2x - 6)$ sit relative to $y = f(2x)$?
  \answerto{Three to the right, not six. Write it as $f\bigl(2(x-3)\bigr)$
  first: the shift is measured after the squash, so it is
  $\frac{6}{2}$.}
\item Two students describe $y = (2x)^{2}$. One says the graph of $x^{2}$
  squashed horizontally by two, the other says stretched vertically by four.
  Which is right?
  \answerto{Both, and the graphs are identical. $(2x)^{2} = 4x^{2}$, so this
  is one of the cases where an inside move and an outside move happen to
  coincide. It is special to squaring and does not generalise.}
\item A unit computes $\sigma(wx + b)$ and you want it to be near $0$ below
  $x = 10$ and near $1$ above it, with the transition taking about one unit of
  $x$. Give $w$ and $b$.
  \answerto{The band from $\num{0.1}$ to $\num{0.9}$ is
  $\frac{\val{f05.band.w1}}{w}$, so $w$ near $\num{4.4}$ gives a band about one
  wide, and the crossing at $10$ then needs $b = -10w$, about $-44$.}
\item Why does raising the weight of a logistic unit without changing the bias
  move its threshold?
  \answerto{Because the threshold is $-\frac{b}{w}$ and the weight is in the
  denominator. Doubling $w$ halves the distance of the threshold from the
  origin, which is a side effect nobody asked for and a good reason to
  parametrise by the threshold instead.}
\item Compose $\relu$ with itself. What do you get?
  \answerto{$\relu$ again: $\max(0, \max(0, x)) = \max(0, x)$. A function that
  is its own composition is called idempotent, and it means that stacking two
  of them buys nothing at all.}
\item $f(x) = x^{2}$ has no inverse over the reals. Give two different
  restrictions of its domain that each give it one, and the inverse in each
  case.
  \answerto{$x \ge 0$ gives $\sqrt{y}$; $x \le 0$ gives $-\sqrt{y}$. The
  inverse is a property of the function \emph{and} its domain, not of the
  formula.}
\item The graph of $f$ passes through $(0, 1)$ and $(1, e)$. Where does the
  graph of $f^{-1}$ pass?
  \answerto{Through $(1, 0)$ and $(e, 1)$. The coordinates swap, which is the
  reflection in $y = x$ done one point at a time.}
\item Is $\sigma$ strictly increasing? Does it therefore preserve the
  $\argmax$ of a list of scores taken one at a time?
  \answerto{Yes to both. A binary classifier's ranking is therefore the same
  whether you sort by logit or by probability, which is why an
  area-under-the-curve figure does not change when you drop the sigmoid.}
\item A loss is reported as $L$ and a colleague reports $100L$. Do the two
  agree about which of two models is better?
  \answerto{Yes. Multiplying by a positive constant is strictly increasing, so
  the $\argmin$ does not move. They will disagree about every sentence
  containing the word \emph{small}.}
\item Softmax is applied twice: $\softmax(\softmax(z))$. Does the $\argmax$
  move?
  \answerto{No, because softmax is built from strictly increasing steps and a
  division by a positive total. The result is much flatter than one
  application and is not a sensible thing to compute, but the top entry is the
  same one.}
\item At which temperature in the table of frame $43$ is the second-placed
  token nearest to the first, and what does that suggest about sampling there?
  \answerto{At $T = \num{2.0}$, where \val{f05.sm.t20.p1} and
  \val{f05.sm.t20.p2} are close. The ranking is unchanged but the draw is
  nearly a coin toss between them, which is what the setting is for.}
\item Adding a constant $c$ to every logit leaves the softmax unchanged.
  Show it, and say which frame of this program it is an instance of.
  \answerto{$\frac{e^{z_j + c}}{\sum_k e^{z_k + c}}
  = \frac{e^{c}e^{z_j}}{e^{c}\sum_k e^{z_k}}$, and the $e^{c}$ cancels. It is
  frame $45$'s observation that dividing every entry by one positive number
  changes no comparison \dash{} and it is why subtracting the maximum before
  exponentiating is free, which Program~\ref{prog:P02} needs.}
\item Give a strictly decreasing function of a probability that is used as a
  loss, and say what its $\argmin$ corresponds to.
  \answerto{$-\ln p$. Its $\argmin$ over the parameters is the $\argmax$ of
  the likelihood, because $\ln$ is strictly increasing and the minus sign
  strictly decreasing.}
\item Why can the four moves of section 3 never turn a straight line into a
  curve?
  \answerto{Each of them adds a constant to, or multiplies by a constant, the
  input or the output. A linear function of a linear function is linear, which
  is frame $32$'s collapse stated the other way round.}
\end{furtherproblems}
